\documentclass[11pt,a4paper]{article}

% ---------------------------------------------------------
% PREAMBLE (stable, warning-free, Overleaf-safe)
% ---------------------------------------------------------
\usepackage[utf8]{inputenc}
\usepackage[T1]{fontenc}
\usepackage{lmodern}
\usepackage{amsmath, amssymb, amsfonts}
\usepackage{bm}
\usepackage{geometry}
\usepackage{graphicx}
\usepackage{hyperref}
\usepackage{cite}
\usepackage{authblk}
\usepackage{physics}
\usepackage{mathtools}

\geometry{margin=2.4cm}

\title{\textbf{Companion LXIX: Spectral Density Functions of Persistence Diagrams in the Relaxation-Driven Cyclic Cosmology}}
\author{Michael Lehmann \\ \texttt{mi.lehmann@gmx.de}}
\date{April 2026}

\begin{document}
\maketitle

\begin{abstract}
Spectral density functions (SDFs) of persistence diagrams provide a frequency-
resolved characterization of topological fluctuations in cosmological fields. In 
weak-lensing analyses, SDFs quantify how the birth and death of connected 
components, loops, and void-like structures contribute across frequency modes. In 
the standard cosmological model, the SDF reflects the nonlinear matter distribution 
and the projection of large-scale structure. In the Relaxation-Driven Cyclic 
Cosmology (RDCC), the scale-dependent evolution of the gravitational potential 
modifies the spectral distribution of topological features in a characteristic way. 
This Companion develops a continuous description of SDFs in RDCC, deriving how the 
relaxation scale influences frequency-resolved topology, spectral coherence, and 
the observational signatures in weak-lensing surveys. The resulting framework 
provides a distinctive probe of the RDCC gravitational field through spectral 
topology.
\end{abstract}
% ---------------------------------------------------------
\section{Introduction}

Persistence diagrams encode the birth and death of topological features in 
weak-lensing convergence fields. While cross-correlation functions quantify 
filtration-coupled topology, spectral density functions provide a frequency-
resolved view of topological fluctuations.

In the standard cosmological model, the SDF reflects the nonlinear matter 
distribution and the projection of large-scale structure. In RDCC, the 
gravitational potential evolves differently. The relaxation mechanism suppresses 
the decay of small-scale modes and enhances the coherence of large-scale modes. 
This modifies the spectral distribution of topological features in a scale-
dependent manner, producing a distinctive signature in SDF analyses.
% ---------------------------------------------------------
\section{Spectral Density Functions of Persistence Diagrams}

Given a time- or filtration-ordered sequence of persistence diagrams 
$\{D(\nu_{i})\}$, the spectral density function is defined as the Fourier transform 
of the Wasserstein-distance fluctuations:

\begin{equation}
    S_{D}(\omega)
    = \left|
        \sum_{i}
        W_{p}(D(\nu_{i}), D^{*})
        e^{-i \omega \nu_{i}}
      \right|^{2}.
\end{equation}

This function measures the frequency content of topological fluctuations.
% ---------------------------------------------------------
\section{RDCC Modification of Persistence Diagrams}

In RDCC, the convergence evolves as
\begin{equation}
    \kappa_{\mathrm{RDCC}}(\ell)
    = \kappa(\ell)
      \left[
        1 - \chi_{\kappa}
        \left(\frac{\ell}{\ell_{\mathrm{rel}}}\right)^{\alpha}
      \right].
\end{equation}

This modifies the birth and death thresholds of topological features:

\begin{align}
    b_{\mathrm{RDCC}} &= b \left[ 1 - \chi_{b} (\ell/\ell_{\mathrm{rel}})^{\alpha} \right], \\
    d_{\mathrm{RDCC}} &= d \left[ 1 - \chi_{d} (\ell/\ell_{\mathrm{rel}})^{\alpha} \right].
\end{align}

Thus the RDCC spectral sequence becomes
\begin{equation}
    W_{p}(D_{\mathrm{RDCC}}(\nu), D^{*}_{\mathrm{RDCC}}).
\end{equation}
% ---------------------------------------------------------
\section{Spectral Density Functions in RDCC}

The RDCC spectral density function is

\begin{equation}
    S_{D,\mathrm{RDCC}}(\omega)
    = S_{D}(\omega)
      \left[
        1 - \chi_{S}
        \left(\frac{\ell}{\ell_{\mathrm{rel}}}\right)^{\alpha}
      \right].
\end{equation}

This modification reflects the relaxation-driven change in frequency-resolved 
topology.
% ---------------------------------------------------------
\section{Spectral Topological Signatures of RDCC}

The relaxation mechanism enhances the coherence of large-scale modes. Spectral 
density functions therefore exhibit:

\begin{itemize}
    \item enhanced low-frequency power (large-scale coherence),
    \item suppressed high-frequency power (small-scale suppression),
    \item a characteristic turnover near $\ell_{\mathrm{rel}}$,
    \item modified spectral slopes across redshift,
    \item altered frequency alignment of persistence lifetimes.
\end{itemize}

These signatures provide a direct probe of the RDCC gravitational field through 
spectral topology.
% ---------------------------------------------------------
\section{Conclusion}

Spectral density functions of persistence diagrams in the Relaxation-Driven Cyclic 
Cosmology reflect the scale-dependent evolution of the gravitational potential and 
the nonlinear matter distribution. The relaxation mechanism modifies the frequency-
resolved topology in a scale-dependent manner, producing distinctive signatures in 
weak-lensing surveys. These effects provide a direct observational window into the 
relaxation scale and the temporal structure of the RDCC gravitational field. The 
present Companion establishes the theoretical foundation for future studies of 
spectral topology, nonlinear structure, and the interplay between light and gravity 
in RDCC.

\bibliographystyle{apsrev4-2}
\nocite{*}
\bibliography{references}


\end{document}
